\documentclass[10pt]{article}
\usepackage[vietnamese]{babel}
\usepackage[letterpaper, top=18mm, bottom=20mm, left=15mm, right=15mm, headheight=14pt]{geometry}
\usepackage{mathtools,amsmath,amssymb}
\usepackage{newtxtext,newtxmath}
\usepackage{xcolor}
\usepackage{tcolorbox}
\usepackage{fancyhdr}

% Định nghĩa bảng màu Stewart Calculus chuẩn xác
\definecolor{stewartcyan}{RGB}{0, 136, 204}
\definecolor{stewartred}{RGB}{204, 34, 41}

% Cấu hình khung định nghĩa Stewart
\newtcolorbox{definitionbox}{%
    colback=white,
    colframe=stewartred,
    arc=0mm,
    boxrule=0.8pt,
    left=10pt,
    right=10pt,
    top=7pt,
    bottom=7pt
}

% Cấu hình tiêu đề đầu trang
\pagestyle{fancy}
\fancyhf{}
\renewcommand{\headrulewidth}{0pt}
\fancyhead[L]{\textbf{\textsf{178}}\qquad\textbf{\textsf{CHƯƠNG 3}}\quad\textsf{Các quy tắc tính đạo hàm}}
\fancyhead[R]{}

\setlength{\parindent}{0pt}
\setlength{\parskip}{5pt}

\begin{document}

\noindent
\begin{minipage}[t]{0.23\textwidth}
    % Cột lề trái: dóng hàng ngang chính xác với khung Quy tắc tổng và hiệu
    \vspace*{34pt}
    \begin{flushleft}
        \small
        Sử dụng ký hiệu dấu phẩy, ta có thể viết Quy tắc tổng và hiệu dưới dạng
        \begin{align*}
            (f + g)' &= f' + g' \\[3pt]
            (f - g)' &= f' - g'
        \end{align*}
    \end{flushleft}
\end{minipage}%
\hfill
\begin{minipage}[t]{0.74\textwidth}
    % Cột nội dung chính
    Quy tắc tiếp theo cho chúng ta biết rằng \textit{đạo hàm của một tổng (hoặc hiệu) các hàm số bằng tổng (hoặc hiệu) các đạo hàm của chúng}.

    \vspace{1pt}
    \begin{definitionbox}
        \textbf{\textsf{\color{stewartred}Quy tắc Tổng và Hiệu}}\hspace{0.8em}Nếu $f$ và $g$ đều khả vi, thì
        \begin{align*}
            \frac{d}{dx} [f(x) + g(x)] &= \frac{d}{dx} f(x) + \frac{d}{dx} g(x) \\[6pt]
            \frac{d}{dx} [f(x) - g(x)] &= \frac{d}{dx} f(x) - \frac{d}{dx} g(x)
        \end{align*}
    \end{definitionbox}

    \vspace{4pt}
    \noindent\textbf{\textsf{\color{stewartred}CHỨNG MINH}}\hspace{0.8em}Để chứng minh Quy tắc tổng, ta đặt $F(x) = f(x) + g(x)$. Khi đó
    \begin{align*}
        F'(x) &= \lim_{h \to 0} \frac{F(x + h) - F(x)}{h} \\[5pt]
        &= \lim_{h \to 0} \frac{[f(x + h) + g(x + h)] - [f(x) + g(x)]}{h} \\[5pt]
        &= \lim_{h \to 0} \left[ \frac{f(x + h) - f(x)}{h} + \frac{g(x + h) - g(x)}{h} \right] \\[5pt]
        &= \lim_{h \to 0} \frac{f(x + h) - f(x)}{h} + \lim_{h \to 0} \frac{g(x + h) - g(x)}{h} \qquad {\color{stewartcyan}\text{(theo Quy tắc giới hạn 1)}} \\[5pt]
        &= f'(x) + g'(x)
    \end{align*}

    Để chứng minh Quy tắc hiệu, ta viết $f - g$ là $f + (-1)g$ rồi áp dụng Quy tắc tổng và Quy tắc nhân với hằng số.\hfill{\color{stewartred}$\blacksquare$}

    \vspace{4pt}
    Quy tắc tổng có thể được mở rộng cho tổng của một số lượng hàm bất kỳ. Chẳng hạn, áp dụng định lý này hai lần, ta được
    \[
        (f + g + h)' = [(f + g) + h]' = (f + g)' + h' = f' + g' + h'
    \]

    Quy tắc nhân với hằng số, Quy tắc tổng và Quy tắc hiệu có thể được kết hợp với Quy tắc lũy thừa để lấy đạo hàm của bất kỳ đa thức nào, như ba ví dụ sau đây sẽ minh họa.

    \vspace{6pt}
    \noindent\textbf{\textsf{\color{stewartcyan}VÍ DỤ 5}}
    \vspace{-4pt}
    \begin{align*}
        \MoveEqLeft[1.5] \frac{d}{dx} (x^8 + 12x^5 - 4x^4 + 10x^3 - 6x + 5) \\[2pt]
        &= \frac{d}{dx}(x^8) + 12\frac{d}{dx}(x^5) - 4\frac{d}{dx}(x^4) + 10\frac{d}{dx}(x^3) - 6\frac{d}{dx}(x) + \frac{d}{dx}(5) \\[4pt]
        &= 8x^7 + 12(5x^4) - 4(4x^3) + 10(3x^2) - 6(1) + 0 \\[4pt]
        &= 8x^7 + 60x^4 - 16x^3 + 30x^2 - 6 \tag*{\color{stewartcyan}\blacksquare}
    \end{align*}
\end{minipage}

\end{document}